% =====================================================================
%  THÈME 3 — Proportions stœchiométriques
%  Niveau MOYEN — Exercices
% =====================================================================
\documentclass[11pt,a4paper]{article}
\usepackage[utf8]{inputenc}
\usepackage[T1]{fontenc}
\usepackage{lmodern}
\usepackage[french]{babel}
\usepackage{amsmath,amssymb,mathtools}
\usepackage{icomma}
\usepackage[version=4]{mhchem}
\usepackage{siunitx}
\sisetup{output-decimal-marker={,},per-mode=symbol}
\DeclareSIUnit\Molar{\mole\per\litre}
\usepackage{xcolor}
\usepackage{tikz}
\usepackage[most]{tcolorbox}
\usepackage{enumitem}
\usepackage{array}
\usepackage{fancyhdr}
\usepackage[a4paper,margin=2.1cm,headheight=26pt,headsep=8pt]{geometry}
\usetikzlibrary{arrows.meta,calc}

\definecolor{primary}{RGB}{142,93,168}
\definecolor{primarydark}{RGB}{82,45,108}
\definecolor{moycol}{RGB}{198,120,9}
\definecolor{reacA}{RGB}{0,90,200}
\definecolor{reacB}{RGB}{210,30,40}
\definecolor{prodC}{RGB}{0,140,90}
\definecolor{prodD}{RGB}{198,93,9}
\newcommand{\nq}[1]{n_{\text{#1}}}
\newcommand{\Vm}{V_{\mathrm m}}

\pagestyle{fancy}\fancyhf{}
\fancyhead[L]{\small\textcolor{primarydark}{\textbf{Fiche 5} — Thème 3 : Proportions}}
\fancyhead[R]{\small\textcolor{moycol}{\textsc{Moyen}}}
\fancyfoot[C]{\small\thepage}
\renewcommand{\headrulewidth}{0.8pt}
\renewcommand{\headrule}{\hbox to\headwidth{\color{primary}\leaders\hrule height \headrulewidth\hfill}}

\newtcolorbox[auto counter]{exo}[1][]{enhanced,breakable,boxrule=1pt,arc=3pt,
  left=3mm,right=3mm,top=2.4mm,bottom=2.4mm,colback=moycol!6,colframe=moycol,
  fonttitle=\bfseries,coltitle=white,
  title={Exercice~\thetcbcounter\ \ifx\relax#1\relax\relax\else\textmd{--- \itshape #1}\fi}}

\setlength{\parindent}{0pt}
\setlength{\parskip}{3pt}

\begin{document}

\begin{center}
{\LARGE\bfseries\color{primarydark}Thème 3 — Proportions}\\[1pt]
{\large\color{moycol}\textbf{Niveau Moyen}}\\[2pt]
{\itshape Objectif : combiner une conversion (masse ou volume) avec la règle de trois. Rendement \SI{100}{\percent}.}\\
{\color{primary}\rule{0.6\linewidth}{1pt}}
\end{center}

\textbf{Données :} $\Vm=\SI{22,4}{\litre\per\mole}$ (TPN).\quad Masses molaires (\si{\gram\per\mole})
fournies dans chaque énoncé.

\begin{exo}[compléter le schéma « en escalier » : synthèse d'un sel]
On veut préparer $\nq{Al2(SO4)3}=\SI{0,300}{\mole}$ de sulfate d'aluminium selon
\[ \ce{2 Al(OH)3 + 3 H2SO4 -> Al2(SO4)3 + 6 H2O}. \]
Reproduis et \textbf{complète} le schéma ci-dessous (cases « ? ») en remontant depuis le produit visé.
\begin{center}
\begin{tikzpicture}[font=\footnotesize,scale=0.95]
  \def\cA{0}\def\cB{3.4}\def\cArr{6.2}\def\cC{8.4}\def\cD{10.9}
  \node[reacA,font=\bfseries] at (\cA,2.3) {2\,\ce{Al(OH)3}};
  \node at (\cA+1.7,2.3) {$+$};
  \node[reacB,font=\bfseries] at (\cB,2.3) {3\,\ce{H2SO4}};
  \draw[->,>=Stealth,line width=1pt] (\cArr-0.7,2.3) -- (\cArr+0.7,2.3);
  \node[prodC,font=\bfseries] at (\cC,2.3) {\ce{Al2(SO4)3}};
  \node at (\cC+1.55,2.3) {$+$};
  \node[prodD,font=\bfseries] at (\cD,2.3) {6\,\ce{H2O}};
  \node[black!60] at (-1.9,1.5) {$t_0$};
  \node[reacA] at (\cA,1.5) {\textbf{?} mol}; \node[reacB] at (\cB,1.5) {\textbf{?} mol};
  \node[prodC] at (\cC,1.5) {$0$}; \node[prodD] at (\cD,1.5) {$0$};
  \node[black!60] at (-1.9,0.4) {$t_f$};
  \node[reacA] at (\cA,0.4) {$0$}; \node[reacB] at (\cB,0.4) {$0$};
  \node[prodC,font=\bfseries] at (\cC,0.4) {$0{,}300$ mol}; \node[prodD] at (\cD,0.4) {\textbf{?} mol};
  \draw[->,>=Stealth,primary,line width=0.9pt] (\cC,0.7) to[bend left=14] node[above,black,font=\scriptsize]{$\times 2$} (\cA,1.2);
  \draw[->,>=Stealth,primary,line width=0.9pt] (\cC,0.7) to[bend left=9] node[above,pos=0.6,black,font=\scriptsize]{$\times 3$} (\cB,1.2);
  \draw[->,>=Stealth,primary,line width=0.9pt] (\cC,0.35) to[bend right=16] node[below,black,font=\scriptsize]{$\times 6$} (\cD,0.35);
  \draw[primary!30,rounded corners=3pt] (-2.5,-0.1) rectangle (12.2,2.7);
\end{tikzpicture}
\end{center}
Donne $\nq{Al(OH)3}$, $\nq{H2SO4}$ et $\nq{H2O}$.
\end{exo}

\begin{exo}[combustion : produire une masse de \ce{CO2} donnée]
On veut produire \SI{88}{\gram} de dioxyde de carbone par la combustion du méthane :
$\ce{CH4 + 2 O2 -> CO2 + 2 H2O}$. \emph{Données :} $M(\ce{CO2})=44$, $M(\ce{CH4})=16$.
\textbf{a)} Combien de moles de \ce{CO2} cela représente-t-il ?\quad
\textbf{b)} Quelle masse de méthane faut-il brûler ?\quad
\textbf{c)} Quel volume de \ce{O2} (aux TPN) est consommé ?
\end{exo}

\begin{exo}[du carbure de calcium à l'acétylène]
Le carbure de calcium réagit avec l'eau selon
$\ce{CaC2 + 2 H2O -> C2H2 + Ca(OH)2}$. On engage \SI{320}{\gram} de \ce{CaC2}
($M=\SI{64}{\gram\per\mole}$). Combien de moles d'acétylène \ce{C2H2} obtient-on ?
\end{exo}

\begin{exo}[quelles masses de réactifs pour un produit donné ?]
On souhaite fabriquer \SI{34}{\gram} d'ammoniac par $\ce{N2 + 3 H2 -> 2 NH3}$.
\emph{Données :} $M(\ce{NH3})=17$, $M(\ce{H2})=2$, $M(\ce{N2})=28$.
\textbf{a)} Combien de moles de \ce{NH3} veut-on ?\quad
\textbf{b)} Quelle masse de \ce{H2} et quelle masse de \ce{N2} faut-il engager (proportions stœchiométriques) ?
\end{exo}

\begin{exo}[un volume de gaz dégagé]
On décompose \SI{25}{\gram} de calcaire : $\ce{CaCO3 -> CaO + CO2}$
($M(\ce{CaCO3})=\SI{100}{\gram\per\mole}$).
Quel volume de \ce{CO2}, mesuré aux TPN, se dégage-t-il ?
\end{exo}

\end{document}
